\section{Introduction}

Retrieval stores are becoming operational memory for agentic AI systems. Unlike static search applications, agents repeatedly insert, update, and query evidence while operating under tail-latency and context-window constraints. A retrieval stack can therefore pass a static relevance benchmark but still fail deployment: newly inserted evidence may not be retrievable in the agent's top-$k$ context, \maxrowp latency may exceed an interaction budget, or a high-quality configuration may impose an unacceptable downstream token footprint. This paper studies whether a lightweight, commodity-hardware benchmark can produce auditable local deployment evidence for these retrieval-memory choices.

Existing evidence is fragmented across static relevance and systems performance. Information-retrieval benchmarks such as BEIR \citep{thakur2021beir} and HotpotQA \citep{yang2018hotpotqa} provide reusable corpora and judgments, but they do not specify adapter conformance, incremental write semantics, context-cost accounting, or deployment objectives. Vector-search and database benchmarks report latency, throughput, and recall surfaces, but those measurements are often separate from task relevance and post-insert retrievability. The resulting gap is a decision problem: an operator needs to know which retrieval configuration is viable under a stated objective, not which engine wins one isolated metric.

\sys addresses this gap by treating benchmarking as a conformance-gated deployment decision audit. A configuration is reportable only after passing a fixed adapter contract. It is then evaluated on workloads that isolate static retrieval, streaming post-insert retrievability, and multi-evidence retrieval, and the final report exposes quality, a normalized context-cost proxy, \maxrowp latency, budget stability, objective sensitivity, and paired decision margins. The protocol is intentionally scoped to local single-node deployment evidence; its purpose is to make the decision assumptions explicit and auditable rather than to rank vector databases universally.

The archived evaluation supports four findings. First, under a strict 200 ms \maxrowp policy, FAISS CPU with \texttt{BAAI/bge-small-en-v1.5} is the policy-selected strict-latency configuration for S1, S2, and S3. Second, this is not a universal ranking: removing the p99 cap changes S3 to LanceDB in-process, and quality-first objectives choose different configurations. Third, paired audits distinguish real margins from fragile aggregate differences: S1 and S2 are tail-latency tie-breaks, and the S3 pgvector-base aggregate edge is indistinguishable from FAISS in a 5{,}000-query matched audit against FAISS exact FlatIP. Fourth, short budget runs are not reliable final-decision substitutes: B0-to-B2 top-1 agreement is absent for S1 and S2, while S3 preserves top-1 despite low full-rank correlation.

The paper contributes:
\begin{enumerate}[leftmargin=*]
  \item A conformance-gated adapter protocol for retrieval infrastructure used as agentic memory, covering lifecycle operations, bulk and incremental writes, query and batch query, payload/vector updates, deletes, flush/commit behavior, index controls, and minimum stats fields.
  \item Three MaxionBench workloads covering static retrieval, streaming post-insert top-10 retrievability, and multi-evidence retrieval under local embedding tiers.
  \item A decision-audit reporting format that separates strict-latency, cost-only/no-p99, and quality-first objectives, then reports quality, normalized context-cost proxy, \maxrowp latency, budget stability, margins, and paired audits.
  \item Empirical local findings showing that strict-latency, cost-only/no-p99, and quality-first objectives select different configurations; that aggregate quality edges can vanish under matched audits; and that short screening runs do not reliably preserve full-budget decisions.
\end{enumerate}

The submitted artifact contains the benchmark code, archived run bundle, validation commands, conformance matrix, generated paper tables and figures, query-level paired-audit summaries, checksums, \hotpotmaxionbench{} metadata, and a claim-evidence map.
